\section{Minerva Contract}

The audit record carries the public inputs the implemented Minerva replay
consumes:

\begin{center}
\small
\begin{tabularx}{\textwidth}{lX}
\toprule
Field & Meaning \\
\midrule
\texttt{method\_id} & Versioned method name: \texttt{minerva-ballot-polling-v1} or \texttt{athena-ballot-polling-v1}. \\
\texttt{risk\_limit\_ppm} & Risk limit in parts per million, used to form the stopping threshold. \\
\texttt{reported\_winner\_votes} & Reported votes for the winner of the two-candidate assertion. \\
\texttt{reported\_loser\_votes} & Reported votes for the loser of the assertion. \\
\texttt{sample\_steps} & Ballot observations. Assorter value \texttt{1/1} is a winner ballot and \texttt{0/1} is a loser ballot. \\
\texttt{round\_index} & Optional zero-based round marker on each sample step. If any step has it, every step must have it and the sequence must be nondecreasing. \\
\texttt{statistic}, \texttt{p\_value\_ppm}, \texttt{decision} & Declared round outputs checked against the replay. \\
\bottomrule
\end{tabularx}
\end{center}

When no \texttt{round\_index} values are present, the run is treated as one
round and the final sample step is the declared round step. When explicit
rounds are present, the final sample step in each round carries the declared
round statistic. For each round, RCOUNT adds that round's observations to the
cumulative sample, counts winner ballots, computes the binomial upper tail
under the reported winner share and under the null share $1/2$, divides those
tails to obtain the likelihood ratio, converts the inverse statistic to a
conservative parts-per-million p-value, and stops when the ratio reaches
$1/\alpha$ for the declared risk limit $\alpha$.
